\chapter{Literature Review}

\section{Classical Yield Curve Models}

The modeling of the term structure of interest rates has long been a central topic in financial
economics. A seminal contribution is the parametric framework proposed by
\textcite{nelson_siegel_1987}, which represents the yield curve using a small number of latent
factors associated with its level, slope, and curvature. The main advantages of this approach
are its parsimony, flexibility, and economic interpretability, which have led to its widespread
use in empirical research and policy institutions.

An important extension was introduced by \textcite{svensson_1994}, who augmented the
original specification with an additional curvature component to improve the fit at longer
maturities. The Nelson--Siegel--Svensson formulation has since become a standard tool for
estimating sovereign yield curves, particularly in central bank applications. Despite their
success in fitting the cross-sectional shape of the yield curve, these static specifications do
not explicitly capture the dynamic evolution of interest rates.

A dynamic interpretation of the Nelson--Siegel framework was provided by
\textcite{diebold_li_2006}, who demonstrated that the latent yield curve factors can be modeled
as time series processes. Their work established a connection between yield curve modeling
and factor-based forecasting methods and laid the foundation for subsequent macro-financial
extensions.

\section{Macroeconomic Influences on the Yield Curve}

A growing body of literature emphasizes the role of macroeconomic fundamentals in shaping
yield curve dynamics. Early empirical evidence suggests that variables such as inflation,
economic activity, and monetary policy expectations are closely related to movements in yields
across maturities. Models that abstract from these relationships may therefore provide an
incomplete description of interest rate behavior.

\textcite{ang_piazzesi_2003} proposed a no-arbitrage framework that integrates
macroeconomic variables into term structure modeling. Their results show that macroeconomic
shocks significantly affect the entire yield curve, highlighting the importance of incorporating
information from the real economy.

Further contributions demonstrate that yield curve factors themselves contain information
about future macroeconomic conditions. \textcite{diebold_rudebusch_aruoba_2006}
developed a macro-finance model that jointly describes macroeconomic variables and yield
curve factors within a unified dynamic system. Their findings reveal strong interactions
between the level and slope of the yield curve and key macroeconomic indicators, reinforcing
the view that financial and macroeconomic dynamics are closely interconnected.

\section{Econometric Approaches to Yield Curve Dynamics}

Econometric methods such as Vector Autoregression (VAR) and Dynamic Factor Models (DFM)
are widely employed to analyze the joint dynamics of yield curve factors and macroeconomic
variables. VAR models allow for flexible dynamic interactions without imposing strong
structural assumptions, making them particularly suitable for empirical macro-finance
applications.

Dynamic factor models, as discussed by \textcite{stock_watson_2002}, provide a parsimonious
representation of information contained in large macroeconomic datasets by summarizing
co-movements through a small number of latent factors. When applied to yield curve modeling,
these approaches offer a natural framework for integrating macroeconomic information into
interest rate dynamics.

Empirical studies indicate that macro-augmented yield curve models often outperform
purely yield-based specifications in forecasting exercises, particularly at medium and long
horizons \parencite{koop_korobilis_2010}. At the same time, these models maintain a relatively
transparent economic interpretation, which is desirable for policy analysis and financial risk
management.

\section{Stochastic Interest Rate Models}

An alternative strand of the literature focuses on stochastic models of interest rates, with
particular emphasis on applications in pricing and risk management. Short-rate models
describe the evolution of an instantaneous interest rate from which the entire term structure
can be derived. Among these, the Hull--White one-factor model remains one of the most widely
used frameworks in practice.

\textcite{hull_white_1990} extended earlier mean-reverting short-rate models by allowing for
time-dependent parameters, enabling an exact fit to the observed initial yield curve. The
analytical tractability of the model and its suitability for pricing interest rate derivatives
have contributed to its widespread adoption in financial institutions. However, the model
abstracts from explicit macroeconomic mechanisms and relies primarily on market-implied
interest rate dynamics.

Due to these characteristics, stochastic short-rate models are commonly used as benchmarks
when evaluating alternative yield curve modeling approaches, particularly those aimed at
improving economic interpretation rather than pricing accuracy.

\section{Summary and Research Gap}

The literature reveals two dominant approaches to yield curve modeling. The first emphasizes
econometric representations of yield curve factors and their relationship with macroeconomic
variables, while the second focuses on stochastic interest rate models designed primarily for
pricing and risk management. Although both strands have been extensively studied, direct
empirical comparisons between macro-augmented econometric yield curve models and
stochastic short-rate models using a consistent dataset and evaluation framework remain
relatively limited.

This thesis addresses this gap by integrating macroeconomic variables into an econometric
yield curve model and benchmarking its performance against the Hull--White stochastic model.
